\section{案例推演：5 类典型任务的端到端执行路径}

\subsection{案例 1：陌生数据格式解析与分析报告生成}

这是访谈中最具代表性的“长尾任务”原型。输入往往是非标准文件格式、零散说明文档和不完整约束。一个稳定的 Agent 执行链路应当是：识别格式 -> 检索解析工具 -> 验证解析正确性 -> 生成中间结构化数据 -> 产出可读报告。与传统脚本不同，通用 Agent 的优势在于可以边执行边修正策略。

\begin{importantbox}{该类任务的成功关键}
\begin{itemize}
    \item 不把第一次解析成功当成结束，必须做一致性校验
    \item 明确区分“解析失败”与“语义理解失败”
    \item 报告输出应附带可追溯中间结果，降低审核成本
\end{itemize}
\end{importantbox}

如果把这一任务放进访谈语境，可以理解为 Manus 的核心能力展示：面对未见过的输入分布，系统并不依赖固定模板，而是通过 Tool Use 和代码执行形成临时解决方案。

\subsection{案例 2：多来源资料整合与决策备忘录}

另一类高频任务是“把多个来源的信息收敛成可决策文本”。这类任务对 Agent 的要求不是“写得长”，而是“信息一致、结构清晰、结论可追溯”。实际执行常见问题包括：来源冲突未标记、事实与观点混写、引用链断裂。

\begin{warningbox}{信息整合任务最容易出现的三类错误}
\begin{enumerate}
    \item 把二手观点当一手事实
    \item 引用存在但不可定位到原始上下文
    \item 结论强度与证据强度不匹配
\end{enumerate}
\end{warningbox}

访谈中强调“用户最终要的是省心决策”，因此这类任务的评测不能只看文本流畅度，还要看结论是否可复核。可以在输出模板中强制加入“证据等级”与“反例提示”字段。

\subsection{案例 3：研究类任务（Deep Research）与可复现交付}

研究类任务常被误解为“多搜一点资料”。实际上，它更接近一个 mini research pipeline：问题拆分、信息检索、证据筛选、假设修正、结论输出。访谈里对 Deep Research 的判断很清楚：它会逐步同质化，所以差异化点在执行稳定性与交付质量。

\begin{knowledgebox}{研究类任务的输出结构建议}
\begin{itemize}
    \item 结论摘要：3--5 条可执行判断
    \item 证据面板：每条判断对应来源、时间、可信度
    \item 不确定性声明：哪些结论需要追加验证
    \item 下一步行动：明确人机分工边界
\end{itemize}
\end{knowledgebox}

如果把这套结构固化到产品中，就能显著降低“看起来很全、实际上不能决策”的输出风险。访谈中的工程哲学正是如此：把模糊能力变成结构化交付。

\subsection{案例 4：工具链编排任务（代码、浏览器、命令行混合）}

这种任务最能体现“统一动作空间”价值。传统流程常把浏览器自动化、脚本执行、结果汇总拆在多个系统里，导致状态断裂。Manus 路线的优势在于把动作统一收敛到会话沙盒中，减少上下文丢失。

\begin{importantbox}{混合工具链任务的执行顺序模板}
\begin{enumerate}
    \item 定义目标状态（最终需要交付什么）
    \item 规划最短可行路径（最少工具、最少步骤）
    \item 执行后立刻校验（而不是最后一次性验收）
    \item 失败时局部回滚（保留已验证的中间成果）
\end{enumerate}
\end{importantbox}

访谈对“删工具”的强调，在这里尤为重要。工具越多，路由错误和状态漂移概率越高。相比增加能力表面宽度，更应提升默认路径的可靠性。

\subsection{案例 5：长任务后台执行与人工接管策略}

长任务是通用 Agent 价值最高也最脆弱的场景。访谈中“沙盒持续运行”解决的是会话连续性问题，但还需要明确人工接管边界。否则任务在异常时要么无限重试，要么静默失败。

\begin{warningbox}{长任务必须显式定义的三条接管规则}
\begin{itemize}
    \item 何时自动重试、重试多少次
    \item 何时必须请求人工确认
    \item 何时触发安全中止并保留现场
\end{itemize}
\end{warningbox}

这类规则一旦缺失，系统会出现“看似自动化，实则不可控”的伪自动化。访谈中的可观测性与回放机制，正是为了解决这个问题。

\subsection{跨案例对照：能力、风险、治理动作}

\begin{center}
\begin{tabular}{p{3.0cm}p{3.5cm}p{4.1cm}p{3.5cm}}
\toprule
任务类型 & 关键能力 & 主要风险 & 治理动作 \\
\midrule
陌生格式解析 & Tool discovery + coding & 解析正确但语义错读 & 双层校验（结构 + 语义） \\
资料整合备忘录 & 证据归并 +结构化写作 & 引用链断裂 & 强制证据字段 \\
研究任务 & 问题拆解 + 证据筛选 & 结论强度过度外推 & 输出不确定性声明 \\
混合工具链 & 状态管理 + 回滚能力 & 工具误路由 & 工具白名单与默认路径 \\
长任务执行 & 持续运行 + 监控告警 & 静默失败或无限重试 & 接管规则与中止条件 \\
\bottomrule
\end{tabular}
\end{center}

\subsection{本章小结}

5 类案例共同指向同一结论：通用 Agent 的价值不在“会不会做”，而在“能不能稳定做完并可追溯地交付”。这也是访谈里反复强调系统工程与组织纪律的根本原因。


